if x <= y then
    z := (y - x) / 2 + x;
else
    z := (x - y) / 2 + y;
z := 2 * z ;
assert z == x + y;

%Calculation of WLP(Trick, true):

x <= y => (2 * ((y-x) / 2 + x) == x + y && true)
&&
! (x <= y) => (2 * ((x-y) / 2 + y) == x + y && true)
    if x <= y then {
        2 * ((y-x) / 2 + x) == x + y && true
        z:= (y-x) / 2 + x;
        2 * z == x+y && true
        z:= 2 * z;
        z == x+y && true
    }
    else {
        2 * ((x-y) / 2 + y) == x + y && true 
        z:= (x-y) / 2 + y;
        2 * z == x + y && true
        z:= 2 * z;
        z == x+y && true
    }
    z == x+y && true
assert (z == x+y);
    true

WLP (Trick, true) is
    x <= y => (2 * ((y-x) / 2 + x) == x + y && true)
                        &&
! (x <= y) => (2 * ((x-y) / 2 + y) == x + y && true)

2
(a)
Para a primeira conjuncao:
; Variable declarations
(declare-const x Int)
(declare-const y Int)

; Constraints
(assert (<= x y))
(assert (not
    (= (* 2 (+ (div (- y x) 2) x))
       (+ x y))))
; Solve
(check-sat)
(get-model)
Resposta: SAT, ou seja, achou um contra exemplo, logo não é válido.
Exemplo: x = 0 e y = 1 -> 2 * ((1-0) / 2+0) != 0 + 1  <=> 2 * (1 / 2 + 0) != 1 <=> 2 * 0 (integer division) != 1
<=> 0 != 1 (verdadeiro) -> contra exemplo achado

Para a segunda conjunção:

(declare-const x Int)
(declare-const y Int)

(assert (> x y))
(assert (not
    (= (* 2 (+ (div (- x y) 2) y))
       (+ x y))))

(check-sat)
Resposta: SAT, ou seja, achou um contra exemplo, logo não é válido.
Exemplo: x = 1 e y = 0 -> 2 * ((1-0) / 2+0) != 0 + 1 <=> 2 * (1 / 2 + 0) != 1 <=> 2 * 0 (integer division) != 1
<=> 0 != 1 (verdadeiro) -> contra exemplo achado

RESPOSTA FINAL: Trick não é correto para todos os valores iniciais.

(b)
(a resposta meio que ja foi dada na alinea anterior. Quando passar para o latex fazer separado).
O solver retorna 'sat' para as duas queries, tendo achado contra exemplos para ambas as conjunções do WLP.
Para a primeira conjunção, um contra exemplo é x=0, y=1. O ramo then produz z=0, que é diferente de x+y=1.
Para a segunda conjunção, um contra exemplo é x=1, y=0. O ramo else produz z=0, que é diferente de x+y=1.
Para o WLP não ser válido, basta que exista, pelo menos, uma conjunção que não seja válida.
Assim, nenhuma das conjunções é válida, pelo que o WLP não é válido para todos os initial values das variáveis.
Podemos concluir que Trick não é correto para todos os integer values de x e y.

Para confirmar, fazendo para a primeira conjunção:
(declare-const x Int)
(declare-const y Int)

(assert (<= x y))
(assert (not
    (= (* 2 (+ (div (- y x) 2) x))
       (+ x y))))
(check-sat)
(get-model)
o resultado obtido é:
sat
(model 
  (define-fun y () Int
    0)
  (define-fun x () Int
    (- 1))
), ou seja, achou como contra-exemplo para a primeira conjunção x=-1 e y=0 -> 2 * ((0-(-1)) / 2 + (-1)) != -1 + 0 <=> -2 != -1, que é verdadeiro.